\documentclass[12pt,a4paper]{article}
\usepackage[utf8]{inputenc}
\usepackage[brazil]{babel}
\usepackage{geometry}
\usepackage{booktabs}
\usepackage{longtable}
\usepackage{hyperref}
\usepackage{xcolor}
\usepackage{fancyhdr}

\geometry{margin=2.5cm}

\pagestyle{fancy}
\fancyhf{}
\rhead{SIVIRA - Sistema de Planejamento de Produção}
\lhead{Modificações de Tempo Máximo de Espera}
\rfoot{Página \thepage}

\title{\textbf{Relatório de Modificações}\\[0.5cm]
\large Tempo Máximo de Espera nas Atividades}
\author{SIVIRA - Sistema de Planejamento de Produção}
\date{11/12/2024}

\begin{document}

\maketitle

\section{Introdução}

Este documento apresenta as modificações realizadas no campo \texttt{tempo\_maximo\_de\_espera}
dos arquivos JSON de atividades. Aproximadamente 50\% das atividades que possuíam tempo
igual a \texttt{00:00:00} foram selecionadas aleatoriamente e receberam novos valores
entre \texttt{00:03:00} e \texttt{00:30:00}.

\section{Resumo das Modificações}

\begin{itemize}
    \item \textbf{Total de modificações em PRODUTOS:} 37
    \item \textbf{Total de modificações em SUBPRODUTOS:} 12
    \item \textbf{Total geral:} 49
\end{itemize}

\section{Modificações em PRODUTOS}

\begin{longtable}{p{4.5cm}p{5.5cm}cc}
\toprule
\textbf{Arquivo} & \textbf{Atividade} & \textbf{Antes} & \textbf{Depois} \\
\midrule
\endhead

\multicolumn{4}{l}{\textbf{1001\_pao\_frances.json}} \\
& pesagem\_de\_massas & 00:00:00 & 00:05:00 \\
& divisao\_de\_massas & 00:00:00 & 00:05:00 \\
& modelagem\_de\_pao & 00:00:00 & 00:20:00 \\
& corte\_de\_pestana & 00:00:00 & 00:05:00 \\
\midrule

\multicolumn{4}{l}{\textbf{1002\_pao\_hamburguer.json}} \\
& pesagem\_de\_massas & 00:00:00 & 00:03:00 \\
& divisao\_de\_massas & 00:00:00 & 00:07:00 \\
& modelagem\_de\_pao & 00:00:00 & 00:15:00 \\
& preparo\_para\_fermentacao & 00:00:00 & 00:10:00 \\
\midrule

\multicolumn{4}{l}{\textbf{1003\_pao\_de\_forma.json}} \\
& pesagem\_de\_massas & 00:00:00 & 00:30:00 \\
& divisao\_de\_massas & 00:00:00 & 00:07:00 \\
& modelagem\_de\_pao & 00:00:00 & 00:03:00 \\
& preparo\_para\_fermentacao & 00:00:00 & 00:20:00 \\
\midrule

\multicolumn{4}{l}{\textbf{1004\_pao\_baguete.json}} \\
& pesagem\_de\_massas & 00:00:00 & 00:10:00 \\
& divisao\_de\_massas & 00:00:00 & 00:15:00 \\
& modelagem\_de\_pao & 00:00:00 & 00:05:00 \\
& corte\_de\_pestana & 00:00:00 & 00:07:00 \\
\midrule

\multicolumn{4}{l}{\textbf{1005\_pao\_tranca\_queijos.json}} \\
& pesagem\_de\_massas & 00:00:00 & 00:20:00 \\
& divisao\_de\_massas & 00:00:00 & 00:03:00 \\
& modelagem\_de\_pao & 00:00:00 & 00:30:00 \\
\midrule

\multicolumn{4}{l}{\textbf{1055\_coxinha\_de\_frango.json}} \\
& montagem\_de\_salgado & 00:00:00 & 00:10:00 \\
& empanamento & 00:00:00 & 00:05:00 \\
\midrule

\multicolumn{4}{l}{\textbf{1059\_folhado\_carne\_sol.json}} \\
& montagem\_de\_salgado & 00:00:00 & 00:15:00 \\
& preparo\_para\_coccao & 00:00:00 & 00:07:00 \\
\midrule

\multicolumn{4}{l}{\textbf{1069\_coxinha\_carne\_sol.json}} \\
& montagem\_de\_salgado & 00:00:00 & 00:03:00 \\
& empanamento & 00:00:00 & 00:20:00 \\
\midrule

\multicolumn{4}{l}{\textbf{1070\_coxinha\_camarao.json}} \\
& montagem\_de\_salgado & 00:00:00 & 00:30:00 \\
& empanamento & 00:00:00 & 00:10:00 \\
\midrule

\multicolumn{4}{l}{\textbf{1071\_coxinha\_queijos.json}} \\
& montagem\_de\_salgado & 00:00:00 & 00:05:00 \\
& empanamento & 00:00:00 & 00:15:00 \\
\midrule

\multicolumn{4}{l}{\textbf{1072\_folhado\_frango.json}} \\
& montagem\_de\_salgado & 00:00:00 & 00:07:00 \\
& preparo\_para\_coccao & 00:00:00 & 00:03:00 \\
\midrule

\multicolumn{4}{l}{\textbf{1073\_folhado\_camarao.json}} \\
& montagem\_de\_salgado & 00:00:00 & 00:20:00 \\
& preparo\_para\_coccao & 00:00:00 & 00:10:00 \\
\midrule

\multicolumn{4}{l}{\textbf{1074\_folhado\_queijos.json}} \\
& montagem\_de\_salgado & 00:00:00 & 00:30:00 \\
& preparo\_para\_coccao & 00:00:00 & 00:05:00 \\

\bottomrule
\end{longtable}

\section{Modificações em SUBPRODUTOS}

\begin{longtable}{p{4.5cm}p{5.5cm}cc}
\toprule
\textbf{Arquivo} & \textbf{Atividade} & \textbf{Antes} & \textbf{Depois} \\
\midrule
\endhead

\multicolumn{4}{l}{\textbf{2002\_massa\_suave.json}} \\
& mistura\_de\_massas & 00:00:00 & 00:10:00 \\
\midrule

\multicolumn{4}{l}{\textbf{2003\_massa\_frituras.json}} \\
& mistura\_de\_massas & 00:00:00 & 00:15:00 \\
\midrule

\multicolumn{4}{l}{\textbf{2004\_massa\_folhados.json}} \\
& laminacao\_1 & 00:00:00 & 00:07:00 \\
& laminacao\_2 & 00:00:00 & 00:20:00 \\
\midrule

\multicolumn{4}{l}{\textbf{2005\_massa\_bolo\_chocolate.json}} \\
& mistura\_de\_massas & 00:00:00 & 00:05:00 \\
\midrule

\multicolumn{4}{l}{\textbf{2008\_carne\_sol\_refogada.json}} \\
& preparo\_armazenamento & 00:00:00 & 00:03:00 \\
& coccao & 00:00:00 & 00:30:00 \\
\midrule

\multicolumn{4}{l}{\textbf{2009\_frango\_refogado.json}} \\
& preparo\_armazenamento & 00:00:00 & 00:10:00 \\
& coccao & 00:00:00 & 00:15:00 \\
\midrule

\multicolumn{4}{l}{\textbf{2010\_creme\_queijo.json}} \\
& preparo\_coccao & 00:00:00 & 00:07:00 \\
& coccao & 00:00:00 & 00:20:00 \\
\midrule

\multicolumn{4}{l}{\textbf{2011\_creme\_camarao.json}} \\
& preparo\_coccao & 00:00:00 & 00:03:00 \\
& coccao & 00:00:00 & 00:30:00 \\
\midrule

\multicolumn{4}{l}{\textbf{2014\_creme\_limao.json}} \\
& preparo\_coccao & 00:00:00 & 00:05:00 \\

\bottomrule
\end{longtable}

\section{Tempos Utilizados}

Os seguintes valores de tempo máximo de espera foram utilizados nas modificações:

\begin{itemize}
    \item \texttt{00:03:00} - 3 minutos
    \item \texttt{00:05:00} - 5 minutos
    \item \texttt{00:07:00} - 7 minutos
    \item \texttt{00:10:00} - 10 minutos
    \item \texttt{00:15:00} - 15 minutos
    \item \texttt{00:20:00} - 20 minutos
    \item \texttt{00:30:00} - 30 minutos (máximo)
\end{itemize}

\section{Observações}

\begin{enumerate}
    \item As atividades foram selecionadas de forma aleatória usando seed 42 para reprodutibilidade.
    \item Atividades que já possuíam tempo diferente de \texttt{00:00:00} foram preservadas.
    \item A modificação do tempo máximo de espera permite flexibilidade no agendamento das atividades.
    \item Os arquivos JSON originais foram modificados diretamente.
\end{enumerate}

\end{document}
